% 本次提交从两个数组分别取一项：CPU 决定工作帧，交换链给出图像编号。
\begin{tikzpicture}[x=1mm,y=1mm,
  every node/.style={font=\sffamily\fontsize{10.2}{13}\selectfont,
    text=Ink,inner sep=0pt,outer sep=0pt},
  head/.style={anchor=west,font=\sffamily\bfseries\fontsize{11.2}{14}\selectfont},
  line/.style={draw=Muted,line width=.65pt,-{Stealth[length=1.9mm]}},
  box/.style={line width=.65pt,rounded corners=1.4mm}]
  \definecolor{FrameBlue}{HTML}{345F84}
  \draw[box,draw=Ink!70] (0,0) rectangle (156,65);
  \draw[box,draw=FrameBlue!80,fill=FrameBlue!6] (4,27) rectangle (75,61);
  \draw[box,draw=Teal!80,fill=Soft] (81,27) rectangle (152,61);
  \node[head,text=FrameBlue] at (7,55.6) {工作帧数组 \normalfont\ttfamily frames};
  \node[head,text=Teal] at (84,55.6) {窗口图像数组 \normalfont\ttfamily images};
  \node[anchor=west] at (7,46.8) {CPU 这次轮到工作帧 0};
  \node[anchor=west] at (84,46.8) {交换链这次交回图像 2};
  \draw[box,draw=FrameBlue,fill=FrameBlue!15] (7,31) rectangle (37,41);
  \draw[box,draw=Rule,fill=white] (40,31) rectangle (70,41);
  \node at (22,36) {帧 0}; \node[text=Muted] at (55,36) {帧 1};
  \foreach \leftEdge/\idx in {84/0,106/1} {
    \draw[box,draw=Rule,fill=white] (\leftEdge,31) rectangle ++(19,10);
    \node[text=Muted] at (\leftEdge+9.5,36) {图 \idx};
  }
  \draw[box,draw=Teal,fill=Teal!16] (128,31) rectangle (147,41);
  \node at (137.5,36) {图 2};
  \draw[line] (22,31)--(22,23)--(51,23)--(51,17.8);
  \draw[line] (137.5,31)--(137.5,23)--(105,23)--(105,17.8);
  \draw[box,draw=Rust!85,fill=Rust!7] (20,4) rectangle (136,17);
  \node at (78,10.5) {用帧 0 的命令，将工具栏和场景面板画到图 2。};
\end{tikzpicture}
